\begin{figure}[t]
    \centering 
    \resizebox{\columnwidth}{!}{%
    \includegraphics{img/comparison.pdf}
    }
    \caption{\textbf{Comparison of different paradigms for learning and information seeking.} \system enables humans to observe and participate in a collaborative discourse among LM agents with different roles. Users can request the system to generate a full-length cited report based on the discourse history and the information collected.}
    \label{fig:compare_paradigm}
\vspace{-1.5em}
\end{figure}

Recent advancements in language models (LMs) \citep{bai2022constitutional, gpt4technicalreport, geminiteam2024gemini} and retrieval-augmented generation (RAG)~\citep{lewis2021retrievalaugmented} have led to more capable chatbots and emerging generative search engines~\citep{liu-etal-2023-evaluating}. %,  significantly transforming how people seek and access information~\citep{liu-etal-2023-evaluating}. 
Compared to traditional search engines and information retrieval (IR) models~\citep{Robertson1977THEORIESAM}, these systems fulfill user queries by generating direct responses, effectively addressing \textit{known unknowns}, where users are aware of their information needs. 

However, a gap remains in using these systems for complex information-seeking scenarios, such as academic research, market analysis, and decision-making, where the system should expose users to their \textit{unknown unknowns} to facilitate knowledge discovery. While the concept of \textit{unknown unknowns} originally referred to unexpected risks in the military, it is linked to the serendipitous discovery of information in the information research context~\citep{foster2003serendipity,agarwal2015towards}. Specifically, %in the field of economics, 
\citet{kirzner1997entrepreneurial} directly contrasts such discovery (``the realization that one had overlooked something in fact readily available'') with successful search (``the deliberate production of information which one knew one had lacked'').

%\citet{Bates1989TheDO} characterizes complex information seeking with the ``berrypicking'' model, where queries evolve dynamically towards a general information need, and the process should lead to discovery and learning rather than a best retrieved set. This dynamic and exploratory nature makes designing supportive systems challenging. Traditional search engines and RAG-based chatbots require users to lead the information-seeking process with search queries or conversational questions, often inducing echo chamber effects \citep{sharma2024generative}. Moreover, users with limited prior knowledge may even struggle to come up with questions~\citep{kuhlthau1991inside,belkin1982ask}.

Prior work on automated expository writing attempts to help readers  navigate the terrain of \textit{unknown unknowns} by curating information from various sources into unified articles with substantial breadth and depth~\citep{shen2023summarization}. %, similar to well-edited Wikipedia pages or textbooks. 
In particular, the STORM writing system demonstrates that LMs paired with search engines can automatically generate a high-quality draft of Wikipedia-like articles on arbitrary topics~\citep{shao2024assisting}. However, producing the static report as the final outcome, STORM does not support any user interaction which is crucial in complex information seeking %\citet{Bates1989TheDO} characterizes this aspect with the ``berrypicking'' model, 
where there is no single, gold query, but queries evolve dynamically towards a goal~\citep{Bates1989TheDO}. This dynamic and exploratory nature makes designing assistance systems challenging. Traditional search engines and RAG chatbots passively react to users' search queries or conversational questions, often inducing echo chamber effects \citep{sharma2024generative} or high cognition load as users with limited prior knowledge may even struggle to formulate questions~\citep{kuhlthau1991inside,belkin1982ask}.


%STORM instructs LMs to ask insightful questions by prompting them with diverse perspectives and simulating conversations, and finally synthesizes collected information into articles with hierarchical outlines. However, the curated articles, while informative, insufficiently adapt to users' evolving intents during the information-seeking process. To balance user agency with intelligent scaffolding, we extend STORM---instead of having LMs ask and answer questions turn by turn, we simulate multi-agent collaborative discourses where users could also participate in the discussion at any time to steer the discourse. We refer to this new system for complex information-seeking assistance as \textbf{\system}.

%We aim to bridge the gap in complex information seeking by designing a collaborative discourse system that enables accurate and grounded information retrieval. This system allows both the user and the system to co-steer the information-seeking direction according to the user’s latent intent, organize collected information on the fly, and synthesize it into long-form reports. We decompose the design of this system into three iterative components: (1) \textbf{Grounded Question Answering}: This module retrieves and synthesizes grounded answers for given questions, ensuring that the information is accurate and reliable. (2) \textbf{Discourse Mind Map Organization}: This module dynamically organizes the collected information based on the user’s evolving intent over the course of the discourse, adjusting to changes in the user’s focus and needs. (3) \textbf{Grounded Question Generation}: This module automatically steers the discourse to align with the user’s intent, generating new questions based on the mind map organization. This system decomposition mimics the iterative process of the foraging loop and sensemaking loop \cite{pirolli2005sensemaking}, facilitating a continuous cycle of information seeking, synthesis, and understanding.

To surface \textit{unknown unknowns} and better support user interaction, we propose \textbf{Collaborative STORM} (\textbf{\system}), an information-seeking assistance system that supports collaborative discourse among users and multiple LM agents. Unlike the one-question-one-answer mode of interaction, \system allows users to learn by observing and participating occasionally in the discourse, emulating a common educational scenario~\citep{nussbaum2008collaborative}. 
To facilitate a thought-provoking discourse and serendipitous discovery, \system simulates two agent types grounded in the search engine: \textit{\experts} who participate by asking or answering questions with different perspectives and a \textit{\moderator}, a non-expert who knows enough to ask good questions and steers the discourse. The user can jump in at any time to steer the discourse and inject questions and opinions according to their interest.  \system maintains a dynamic, hierarchical \textit{mind map} to ensure they can easily follow and engage \citep{buzan1974noting}. 
%\TODO{Is Mind map new or old? do you need to cite?} \yucheng{I add citation that to best of knowledge the first one mentioned using mind map for note tagoking to help recall and critical thinking.} 
Upon the conclusion of the discourse, users can request the system to generate a cited report based on the mind map. 

%with different perspectives who can offer  as well as raising questions, and a who manages the discourse by introducing new questions when it becomes repetitive or strays from the user's goal.% Both roles are grounded---the \expert generates utterances from retrieved sources, while the \moderator poses questions from untapped information to shift discussion threads.

%The design of \system involves three major challenges: (1) simulating thought-provoking collaborative discourse to help users find ``unknown unknowns''; (2) grounding the simulated discourse in reliable sources; (3) facilitating effective user engagement. Inspired by studies showing that participants with diverse viewpoints foster more critical discourses~\citep{nussbaum2008collaborative} and that multi-party dialogues require discourse management~\citep{traum2003issues}, we simulate two role types: \textit{\experts} personified with different perspectives, and a \textit{\moderator} managing the discourse. Both roles are grounded---the \expert generates utterances from retrieved sources, while the \moderator poses questions from untapped information to shift discussion threads. By design, the user acts as a third role, able to interject at any point. To ensure the user can easily follow the discourse, \system maintains a dynamic, hierarchical \textit{mind map} that reorganizes as the user guides the discourse. Finally, when the user is satisfied with the information uncovered through the discourse, \system generates a long-form report based on the mind map, discourse history, and the user’s intent.


For evaluation, we introduce WildSeek, a dataset of topics and user goals from real users engaged in complex information-seeking across multiple domains. We propose automatic metrics to assess both discourse trace and final report quality. Our results show that \system significantly outperforms RAG chatbots in surfacing in-depth and serendipitous information while providing a more engaging learning experience.% Moreover, the \moderator role is crucial for the quality of the emulated collaborative discourse.


We further conduct a human evaluation by inviting 20 users with diverse backgrounds to compare \system with a search engine and a RAG chatbot. 70\% preferred \system over the search engine, and 78\% preferred it over the RAG chatbot for the overall information-seeking experience. Participants find that \system facilitates serendipitous discovery and requires less mental effort.%\TODO{unclear: is this positive or negative, the first phrase suggests negative, the rest suggests positive} our detailed evaluation reveals the need to dynamically adjust to users' mental states across different learning/information-seeking stages.


%Human evaluations involve qualitative feedback from users regarding their experience, the system’s ability to align with their evolving intents, and the overall utility of the generated information. A group of users was invited to interact with the collaborative discourse demo for evaluation. They found that ...

Our main contributions include:
\begin{itemize}
%\vspace{-0.5em}
    \item {We propose \system, a novel system that combines collaborative discourse emulation, human interaction, and information organization to assist learning and complex information seeking.}% and build a live demo, showcasing a novel auto-steering collaborative discourse system. This system allows users to dynamically and iteratively forage information and facilitate sense-making.
 %   \vspace{-0.8em}
    \item {We construct the \data dataset from real-world human information-seeking records to evaluate information-seeking assistance tools.}
%  \vspace{-1.6em}
    \item {Results from both automatic and human evaluation show that \system can help humans discover \textit{unknown unknowns} with less mental effort required.}
\end{itemize}